\documentclass[10pt,a4paper]{amsart}
\usepackage{amsmath,amssymb,amsthm,graphicx}
\usepackage{verbatim}
\newenvironment{reference}{\comment}{\endcomment}
%\newenvironment{reference}{}{}
\newenvironment{slogan}{\comment}{\endcomment}
\newenvironment{history}{\comment}{\endcomment}

% For commutative diagrams we use Xy-pic
\usepackage[all]{xy}

% We use 2cell for 2-commutative diagrams.
\xyoption{2cell}
% Theorem environments.
%
\theoremstyle{plain}
\newtheorem{theorem}[subsection]{Theorem}
\newtheorem{proposition}[subsection]{Proposition}
\newtheorem{lemma}[subsection]{Lemma}

\theoremstyle{definition}
\newtheorem{definition}[subsection]{Definition}
\newtheorem{example}[subsection]{Example}
\newtheorem{exercise}[subsection]{Exercise}
\newtheorem{situation}[subsection]{Situation}

\theoremstyle{remark}
\newtheorem{remark}[subsection]{Remark}
\newtheorem{remarks}[subsection]{Remarks}

\numberwithin{equation}{subsection}

% Macros
%
\def\lim{\mathop{\mathrm{lim}}\nolimits}
\def\colim{\mathop{\mathrm{colim}}\nolimits}
\def\Spec{\mathop{\mathrm{Spec}}}
\def\Hom{\mathop{\mathrm{Hom}}\nolimits}
\def\Ext{\mathop{\mathrm{Ext}}\nolimits}
\def\SheafHom{\mathop{\mathcal{H}\!\mathit{om}}\nolimits}
\def\SheafExt{\mathop{\mathcal{E}\!\mathit{xt}}\nolimits}
\def\Sch{\mathit{Sch}}
\def\Mor{\mathop{\mathrm{Mor}}\nolimits}
\def\Ob{\mathop{\mathrm{Ob}}\nolimits}
\def\Sh{\mathop{\mathit{Sh}}\nolimits}
\def\NL{\mathop{N\!L}\nolimits}
\def\CH{\mathop{\mathrm{CH}}\nolimits}
\def\proetale{{pro\text{-}\acute{e}tale}}
\def\etale{{\acute{e}tale}}
\def\QCoh{\mathit{QCoh}}
\def\Ker{\mathop{\mathrm{Ker}}}
\def\Im{\mathop{\mathrm{Im}}}
\def\Coker{\mathop{\mathrm{Coker}}}
\def\Coim{\mathop{\mathrm{Coim}}}

% Boxtimes
%
\DeclareMathSymbol{\boxtimes}{\mathbin}{AMSa}{"02}

%
% Macros for moduli stacks/spaces
%
\def\QCohstack{\mathcal{QC}\!\mathit{oh}}
\def\Cohstack{\mathcal{C}\!\mathit{oh}}
\def\Spacesstack{\mathcal{S}\!\mathit{paces}}
\def\Quotfunctor{\mathrm{Quot}}
\def\Hilbfunctor{\mathrm{Hilb}}
\def\Curvesstack{\mathcal{C}\!\mathit{urves}}
\def\Polarizedstack{\mathcal{P}\!\mathit{olarized}}
\def\Complexesstack{\mathcal{C}\!\mathit{omplexes}}
% \Pic is the operator that assigns to X its picard group, usage \Pic(X)
% \Picardstack_{X/B} denotes the Picard stack of X over B
% \Picardfunctor_{X/B} denotes the Picard functor of X over B
\def\Pic{\mathop{\mathrm{Pic}}\nolimits}
\def\Picardstack{\mathcal{P}\!\mathit{ic}}
\def\Picardfunctor{\mathrm{Pic}}
\def\Deformationcategory{\mathcal{D}\!\mathit{ef}}
\usepackage{cleveref}
\crefname{originallemma}{original lemma}{original lemmas}
\Crefname{originallemma}{Original lemma}{Original lemmas}
\makeatletter
\newlabel{lemma-sheaf-with-exact-support-adjoint}{{0A76}{source}{}{sourcepointer-1}{}}
\newlabel{lemma-sheaf-with-exact-support-adjoint@cref}{{[originallemma][0][]0A76}{[1][0][]source}{}{}{}}
\makeatother
\crefname{originallemma}{original lemma}{original lemmas}
\Crefname{originallemma}{Original lemma}{Original lemmas}
\makeatletter
\newlabel{lemma-sheaf-with-exact-support-ext}{{0A77}{source}{}{sourcepointer-2}{}}
\newlabel{lemma-sheaf-with-exact-support-ext@cref}{{[originallemma][0][]0A77}{[1][0][]source}{}{}{}}
\makeatother
\crefname{originallemma}{original lemma}{original lemmas}
\Crefname{originallemma}{Original lemma}{Original lemmas}
\makeatletter
\newlabel{lemma-sheaf-with-exact-support-quasi-coherent}{{0A78}{source}{}{sourcepointer-3}{}}
\newlabel{lemma-sheaf-with-exact-support-quasi-coherent@cref}{{[originallemma][0][]0A78}{[1][0][]source}{}{}{}}
\makeatother
\crefname{originallemma}{original lemma}{original lemmas}
\Crefname{originallemma}{Original lemma}{Original lemmas}
\makeatletter
\newlabel{more-algebra-lemma-internal-hom-evaluate-tensor-isomorphism}{{0ATK}{source}{}{sourcepointer-4}{}}
\newlabel{more-algebra-lemma-internal-hom-evaluate-tensor-isomorphism@cref}{{[originallemma][0][]0ATK}{[1][0][]source}{}{}{}}
\makeatother
\crefname{originallemma}{original lemma}{original lemmas}
\Crefname{originallemma}{Original lemma}{Original lemmas}
\makeatletter
\newlabel{lemma-flat-bc-sheaf-with-exact-support}{{0E2M}{source}{}{sourcepointer-5}{}}
\newlabel{lemma-flat-bc-sheaf-with-exact-support@cref}{{[originallemma][0][]0E2M}{[1][0][]source}{}{}{}}
\makeatother
\title{The sheaf with exact support of a pseudo-coherent closed immersion commutes with the derived tensor product}
\author{The Stacks Project Authors}
\date{stacks-project:duality.tex:local-snapshot-20261012; GNU FDL 1.2 or later}
\begin{document}
\maketitle
\noindent\textit{Editorial note.} All delivered mathematical text is the author's own, copied from the pinned Stacks Project chapter file duality.tex; no statement or proof step was rewritten, reordered or removed.\par\medskip
\noindent\textit{Editorial note.} Preamble provenance: the pinned chapter file opens with an input of the shared preamble and carries no preamble text of its own, so the project's single Stacks-wide preamble staged at sources/stacks-duality/preamble.tex is excerpted verbatim under the recipe fixed by LANE\_BRIEF 22.1, as three spans, lines 8-13 (verbatim package and the reference/slogan/history environments), lines 14-18 (Xy-pic setup) and lines 154-217 (the native theorem declarations and mathematical macros). Nothing was added, removed or redefined; the wrapper class and the amsmath/amssymb/amsthm/graphicx base packages are supplied by the assembler.\par\medskip
\noindent\textit{Editorial note.} The counted statement and proof contain no cross reference, no citation and no included file, so the delivered body is self contained.\par\medskip
\noindent\textit{Editorial note.} Canonical location of the counted proposition: Stacks tag 0E2N (sources/stacks-duality/pinned-original-source.tex line 2642).\par\medskip
\noindent\textit{Allowed standard prerequisites (editorial).}\begin{itemize}
\item The adjunction and the computation for the sheaf with exact support. Same chapter, Lemmas labelled lemma-sheaf-with-exact-support-adjoint and lemma-sheaf-with-exact-support-ext, cited at proof lines 2675-2680. Convert the isomorphism into a map given by the counit.
\item Quasi-coherence of the sheaf with exact support and the algebraic tensor evaluation. Same chapter, Lemma labelled lemma-sheaf-with-exact-support-quasi-coherent, cited at proof lines 2686-2688, and More on Algebra, Lemma labelled more-algebra-lemma-internal-hom-evaluate-tensor-isomorphism part (3), cited at proof lines 2688-2690. Give the isomorphism after the affine translation.
\item The stalk level alternative. Same chapter, Lemma labelled lemma-flat-bc-sheaf-with-exact-support, cited at proof lines 2690-2692. Provides the second route.
\end{itemize}\medskip
\ifdefined\hypertarget\hypertarget{sourcepointer-1}{}\fi
\noindent\textit{Source reference.} Original lemma 0A76: sources/stacks-duality/pinned-original-source.tex, source lines 2325--2325. Same-chapter pointer into the author environment lemma at source line 2325; invoked at proof line 2660 of the counted proof. It is retained only as a navigation pointer; every load-bearing step is inside the delivered span.\par\medskip
\ifdefined\hypertarget\hypertarget{sourcepointer-2}{}\fi
\noindent\textit{Source reference.} Original lemma 0A77: sources/stacks-duality/pinned-original-source.tex, source lines 2339--2339. Same-chapter pointer into the author environment lemma at source line 2339; invoked at proof line 2661 of the counted proof. It is retained only as a navigation pointer; every load-bearing step is inside the delivered span.\par\medskip
\ifdefined\hypertarget\hypertarget{sourcepointer-3}{}\fi
\noindent\textit{Source reference.} Original lemma 0A78: sources/stacks-duality/pinned-original-source.tex, source lines 2373--2373. Same-chapter pointer into the author environment lemma at source line 2373; invoked at proof line 2665 of the counted proof. It is retained only as a navigation pointer; every load-bearing step is inside the delivered span.\par\medskip
\ifdefined\hypertarget\hypertarget{sourcepointer-4}{}\fi
\noindent\textit{Source reference.} Original lemma 0ATK: sources/stacks-chow/tags.txt, source lines 13618--13618. Cross-chapter standard prerequisite located in the project tag list staged in the local snapshot; invoked at proof line 2667 of the counted proof. It is retained only as a navigation pointer; every load-bearing step is inside the delivered span.\par\medskip
\ifdefined\hypertarget\hypertarget{sourcepointer-5}{}\fi
\noindent\textit{Source reference.} Original lemma 0E2M: sources/stacks-duality/pinned-original-source.tex, source lines 2595--2595. Same-chapter pointer into the author environment lemma at source line 2595; invoked at proof line 2670 of the counted proof. It is retained only as a navigation pointer; every load-bearing step is inside the delivered span.\par\medskip
\begin{lemma}
\label{lemma-sheaf-with-exact-support-tensor}
Let $i : Z \to X$ be a pseudo-coherent closed immersion of schemes.
Let $M \in D_\QCoh(\mathcal{O}_X)$ locally have tor-amplitude in $[a, \infty)$.
Let $K \in D_\QCoh^+(\mathcal{O}_X)$. Then there is a canonical isomorphism
$$
R\SheafHom(\mathcal{O}_Z, K) \otimes_{\mathcal{O}_Z}^\mathbf{L} Li^*M =
R\SheafHom(\mathcal{O}_Z, K \otimes_{\mathcal{O}_X}^\mathbf{L} M)
$$
in $D(\mathcal{O}_Z)$.
\end{lemma}

\begin{proof}
A map from LHS to RHS is the same thing as a map
$$
Ri_*R\SheafHom(\mathcal{O}_Z, K) \otimes_{\mathcal{O}_X}^\mathbf{L} M
\longrightarrow
K \otimes_{\mathcal{O}_X}^\mathbf{L} M
$$
by Lemmas \ref{lemma-sheaf-with-exact-support-adjoint} and
\ref{lemma-sheaf-with-exact-support-ext}. For this map we take the
counit $Ri_*R\SheafHom(\mathcal{O}_Z, K) \to K$ tensored with $\text{id}_M$.
To see this map is an isomorphism under the hypotheses given,
translate into algebra using
Lemma \ref{lemma-sheaf-with-exact-support-quasi-coherent}
and then for example use More on Algebra, Lemma
\ref{more-algebra-lemma-internal-hom-evaluate-tensor-isomorphism} part (3).
Instead of using Lemma \ref{lemma-sheaf-with-exact-support-quasi-coherent}
you can look at stalks as in the second proof of
Lemma \ref{lemma-flat-bc-sheaf-with-exact-support}.
\end{proof}

\end{document}
